%===================================================================================================
\section{Thesis Structure} %++++++++++++++++++++++++++++++++++++++++++++++++++++++++++++++++++++++++
  \label{Thesis Structure} %++++++++++++++++++++++++++++++++++++++++++++++++++++++++++++++++++++++++
%===================================================================================================


\todo


\begin{justify}
    The structure of this thesis is as follows:

    \paragraph{\nameref{Background and Related Work}}
        reviews research on the difficulties novice programmers face when interpreting compiler and test diagnostics,
        and summarises prior approaches that aim to improve the structure and usefulness of error feedback. Building on these perspectives,
        the chapter motivates the thesis goal of presenting failures in a more consistent and
        cognitively simple format.

    \paragraph{\nameref{Problem Analysis}}
        examines typical compiler and testing-framework outputs as they appear in beginner exercises and characterises their technical properties,
        such as verbosity, inconsistency, and the mixing of essential information with framework internals. Based on these observations,
        the chapter derives design requirements for improved test failure reporting, with a stable \emph{tested expression/expected/actual} summary as the primary user-facing interface.

    \paragraph{\nameref{Assertify (and Checkify): Design and Implementation}}
        presents the design rationale and implementation of \emph{Assertify} as a lightweight test and assertion library that wraps existing testing infrastructures.
        It explains how a consistent \texttt{expected=actual} output format is produced, including the use of Quotations and Unquote for expression capture and evaluation,
        as well as wrappers around \texttt{Microsoft.VisualStudio.TestTools.UnitTesting} and \texttt{FsCheck}.
        The chapter also discusses implementation choices and their implications for extensibility.

    \paragraph{\nameref{Results and Discussion}}
        reports the results of the implemented approach and discusses them in relation to the derived requirements and research questions.
        The discussion focuses primarily on technical and qualitative properties achieved by the artifact (e.g., output consistency, coverage of supported assertion forms,
        and integration with unit and property-based testing workflows) and uses illustrative comparisons to baseline framework output to
        highlight differences in interpretability, while separating concrete observations from generalisable claims.

    \paragraph{\nameref{Conclusion}}
        concludes the thesis by summarising the contributions and revisiting the research questions in light of the presented results.
        It then outlines future work, including potential extensions within the wrapper-based approach (e.g., broader expression coverage, richer formatting layers,
        and tighter tooling integration) and longer-term directions that would require more substantial redesign beyond wrappers.
\end{justify}


%===============================================================================
% OLD STRUCTURE +++++++++++++++++++++++++++++++++++++++++++++++++++++++++++++++=
%===============================================================================


\begin{comment}
    This thesis is structured as follows:
    \begin{itemize}
        \item \textbf{Chapter 2: Related Work}
            \skipN Reviews research on novice programmers' difficulties with compiler and test diagnostics,
            and summarises prior approaches to improving the structure and usefulness of error feedback.
            The chapter establishes the conceptual motivation for presenting failures in a more consistent and cognitively simple format.

        \item \textbf{Chapter 3: Analysis of Previous Semester Data}
            \skipN Analyses artefacts from earlier course iterations (e.g., submissions, test outcomes, and common failure patterns)
            to identify recurring sources of errors and to contextualise which kinds of feedback are most frequently encountered by students.
            This chapter motivates the design goals of a novice-oriented reporting scheme by grounding the problem in a concrete teaching setting.

        \item \textbf{Chapter 4: Problem Analysis}
            \skipN Examines typical compiler and testing-framework outputs as they appear in beginner exercises and characterises their technical properties
            (e.g., verbosity, inconsistency, mixing of essential information with framework internals). Based on this analysis, the chapter derives design requirements
            for improved test failure reporting, focusing on a stable \emph{tested expression/expected/actual} summary as the primary interface.

        \item \textbf{Chapter 5: Assertify - Design and Implementation}
            \skipN Presents the design rationale and implementation of \emph{Assertify} as a lightweight test and assertion library that wraps existing testing infrastructures.
            It explains how a consistent \texttt{expected=actual} output format is produced, including the use of Quotations and Unquote for expression capture and evaluation,
            as well as wrappers around \texttt{UnitTesting} and \texttt{FsCheck}. The chapter also discusses implementation choices and their implications for extensibility.

        \item \textbf{Chapter 6: Evaluation and Results}
            \skipN Evaluates Assertify primarily from a technical and qualitative perspective, focusing on the properties the implementation achieves
            (e.g., consistency of output, coverage of supported assertion forms, integration with unit and property-based testing workflows).
            Where applicable, the chapter contrasts Assertify output with baseline framework output to illustrate differences in interpretability,
            while clearly separating observations from generalisable claims.

        \item \textbf{Chapter 7: Discussion and Future Work}
            \skipN Reflects on the contribution and limitations of Assertify and outlines potential extensions.
            This includes incremental improvements within the wrapper approach (e.g., broader expression coverage, richer formatting layers, IDE integration),
            as well as a discussion of what would become possible with more extensive redesign efforts, such as implementing a dedicated testing framework from scratch instead of relying on wrappers.
            The chapter concludes by positioning Assertify as a step toward more pedagogically aligned testing feedback and summarising open directions for future work.

        \item \textbf{Chapter 8: Conclusion}
            \skipN Summarises the thesis contributions, revisits the research questions, and highlights the main insights regarding the feasibility and
            design space of novice-oriented test feedback.
    \end{itemize}

    \skipS This structure emphasizes the design and implementation contribution of the thesis and concludes with a focused discussion of technical possibilities and future extensions,
    ranging from incremental refinements of the current library to more fundamental redesigns of testing feedback infrastructures.
\end{comment}


%===============================================================================
% TODO +++++++++++++++++++++++++++++++++++++++++++++++++++++++++++++++++++++++++
%===============================================================================


% TODO


%===================================================================================================
% EOF ++++++++++++++++++++++++++++++++++++++++++++++++++++++++++++++++++++++++++++++++++++++++++++++
%===================================================================================================